\section{Project Maven: el modelo de detección es solo la entrada de la Decision Chain}

El caso de Project Maven visto en clase parte de la detección de objetivos en datos de sensores que cubren áreas extensas. Aquí no se desarrollan operaciones de ataque; el enfoque está en el problema de sistema: cuando el objeto de interés ocupa una fracción mínima de un campo de visión enorme, la búsqueda manual cuadro por cuadro no es lo bastante rápida; la computer vision puede acortar la fase de observe, pero no puede completar por sí misma la confirmación de identidad, la revisión legal, la asignación de recursos, la autorización de la acción ni el análisis posterior de resultados. Un detector de alta precisión que no se integra en un workflow controlado solo genera más alertas que nadie atiende.

\subsection{OODA loop y latency budget}

\term{OODA} significa Observe, Orient, Decide, Act. Observe recopila el estado de los sensores y del entorno; Orient sitúa las observaciones en el contexto de la misión, los antecedentes y las reglas; en Decide, un sujeto con autoridad elige la acción; Act la ejecuta y produce nuevos resultados. La velocidad del ciclo importa, pero cada fase requiere exactitud y control de permisos. El latency budget debe repartirse en todo el loop, no limitarse a optimizar los inference milliseconds.

\begin{figure}[H]
\centering
\includegraphics[width=0.94\textwidth]{images/08-ooda-decision-chain.png}
\caption{Perspectiva de sistema de Project Maven: la detección acelera Observe; el valor proviene de la OODA decision chain completa.}
\figsource{Redibujado a partir de los subtítulos 21:27--27:00, `images/08-ooda-decision-chain.png`.}
\end{figure}

\begin{knowledgebox}{Lectura de la figura: cada flecha requiere provenance y authority}
La salida del detector debe conservar el sensor, la hora, la versión del modelo y la confidence; la fase Orient necesita relacionar otra información de inteligencia y la incertidumbre; la fase Decide debe mostrar quién tiene la autoridad y qué reglas aplican; después de Act deben registrarse el resultado y los efectos secundarios. Sin estos vínculos, los falsos positivos del modelo no pueden rastrearse y el juicio humano tampoco puede aprovecharse para mejorar la siguiente iteración.
\end{knowledgebox}

\begin{importantbox}{Términos clave: Decision chain y Ontology}
La decision chain es el proceso organizacional completo que va de la observación a la acción y de ahí a la retroalimentación. Ontology no es aquí un término filosófico, sino un modelo compartido del negocio o de la misión: qué entities existen, cómo se relacionan, qué actions pueden ejecutarse, qué estados y qué permisos tienen. La Ontology hace que sensores, modelos, personas y workflows usen la misma capa semántica, y evita que cada sistema describa el mismo objeto con sus propios ID y estructuras de tablas.
\end{importantbox}

\subsection{Del hallazgo de objetos a la acción controlada}

El ponente describió cómo, a partir de «encontrar una gran cantidad de objetos candidatos», se siguieron planteando preguntas: cómo confirmar la identidad, cómo elegir los recursos, cuándo se necesita una revisión legal o de políticas, cómo se retroalimentan los resultados. Esta ampliación pone de manifiesto un patrón frecuente en el AI deployment: el primer modelo resuelve una tarea local y visible, pero lo que en realidad bloquea el valor suelen ser los procesos organizacionales posteriores y las rupturas en los datos.

\begin{warningbox}{Límites de la automatización en escenarios de alto riesgo}
La detección, la priorización y las recomendaciones pueden aumentar la capacidad de procesamiento de las personas, pero eso no equivale a autorizar al sistema a tomar por sí mismo decisiones irreversibles. Los permisos, la human review, la verificación por dos personas, los umbrales de confianza y los mecanismos de apelación o corrección deben diseñarse en función de las consecuencias. La accuracy del modelo no sustituye la legalidad ni elimina al sujeto responsable.
\end{warningbox}

\subsection{Resumen de la sección}

El valor didáctico de Maven no está en un modelo de visión concreto, sino en situar la inference dentro del OODA loop. La observación, la orientación, la decisión, la acción y la retroalimentación requieren una semántica, unos permisos y una auditoría unificados. Este esquema también puede aplicarse a empresas comerciales; solo cambian los objetivos, la autorización legal y las consecuencias de un fallo.
